\chapter{再現条件と補足事項}\label{app:repro}
\section{主比較の条件と根拠の対応}
本稿は元の大容量音響データと学習済みモデルを複製せず，本文の図表を得た小さい評価表，保存予測への参照，出典SHA-256を保持する．研究用の再学習・推論と，本稿をPDFへ変換する処理は別である．草稿の再組版にはGPUを必要としない．

\begin{longtable}{L{35mm}L{99mm}}
\caption{主比較の再現に必要な保存条件．}\label{tab:repro}\\
\toprule 条件 & 値・範囲 \\ \midrule\endfirsthead
\toprule 条件 & 値・範囲 \\ \midrule\endhead
対象 & 2025年6月11日・6月18日，36 WAV，1秒，3\,kHz，選別なし \\
外側分割 & 各WAV内45学習・15評価，1,620/540区間，test fraction 0.25，split seed 42 \\
内部検証 & 学習側chunk 3-fold，seed 42，OOF予測で重みfit \\
学習方針 & clean\_only，モデル・PCA・scaler・重みを全7条件へ固定 \\
CNN学習 & 200 epochs．CNN＋Transformer：lr 0.001，batch 12．AlexNet：lr 0.003，batch 8 \\
RF & PCA 100次元，XGBRF，木数100，最大深さ12，subsample 0.6，colsample\_bynode 0.6 \\
HGB & 34特徴，150反復，lr 0.05，最大葉31，最小葉標本20，L2 1.0，early stoppingなし \\
ET & 34特徴，256本，最小葉標本2，max\_features 1.0，bootstrapなし \\
統合制約 & 元3内部比率固定，元3総量25\%以上，追加各5\%以上，clean OOF全域MSE \\
雑音 & 共通基準RMS $0.0004822199855309985$，水流音1系列，基準SNR 0から$-20$\,dBの6段階 \\
閾値 & 日別221.5051102および271.6776816\,\kw{}．近傍は各値の$\pm10\%$ \\
保存単位 & 元予測\wm{}，主表\kw{}．異なる単位のCSVを加算しない \\
出典 & 実行143609，run hash \texttt{da540f2d}，fit \texttt{7560205f8feb} \\
\bottomrule
\end{longtable}

本文の3-MSE，3-perf，3＋ET，5-modelは，保存表の方式キーと一対一に対応する．\path{data/source_manifest.json}には元CSVとコードのhashを，\path{data/}には本文に使った小さい保存表のコピーを置く．10月14日の誤差分解，10月7日の入力対照を含めて，異なる実行が独立反復であると数えない．

\need{D06}{主実行のPython，TensorFlow，scikit-learn，XGBoost，CUDA等の正確な版と実行環境を保存記録で確認し，固定した環境定義として追記する．本稿を組版した環境の版を，過去の学習時の環境として代用しない．}

\section{先行検討と現在の主条件の関係}
以前の検討では，別日への適用，ピーク高さに基づく学習選別，帯域上限の比較，条件別雑音学習および統合方式の比較を行っている．これらは現在の主比較条件を選ぶ経緯として重要であるが，データ分割，学習条件，モデル構成が同一とは限らない．主系列へ数値を混ぜる前に条件を明記する必要がある．

\begin{longtable}{L{31mm}L{48mm}L{50mm}}
\caption{過去の検討を本文の論理に置く場所．}\\
\toprule 検討 & 現在の位置づけ & 本稿への反映 \\ \midrule\endfirsthead
\toprule 検討 & 現在の位置づけ & 本稿への反映 \\ \midrule\endhead
元WAV分離・別日 & 同一録音内評価と異なる一般化の問い & 第2・3章で現在の分割の限界を明記．旧別日値を現行5の成績にしない \\
2.1--2.5\,kHz選別 & 弱音のONB以上区間を除くことが近傍を悪化させ得る検討 & 現行は選別なし．静かな区間を誤ラベルとはしない \\
3/22\,kHz比較 & 周波数上限と学習条件で性能・誤りが変わる補助比較 & 3\,kHzが物理的に唯一の有効帯域という結論は採らない \\
旧雑音生成診断 & 元音強度と付加音強度の情報経路の確認 & 第3章で固定基準を定義．修正とclean\_only採用を分ける \\
追加モデル・特徴 & 異なる表現・学習器による予測の比較 & 第3章の理由，第4章の性能，第5章の依存へ配置 \\
説明性の数値確認 & 可視化が出力できる状態と，収束・解釈の妥当性の区別 & 中心根拠は共通入力対照と保存予測．未確認画像へ強い主張を依存させない \\
\bottomrule
\end{longtable}

この整理は，過去の作業を再実行する指示ではない．完了している検算と対照は既存の証拠として利用し，現在の結論が必要とする不足だけを次の検討対象にする．

\section{IGを補助資料として扱う際の条件}
回帰関数$F$，入力$x$，基準入力$x'$に対し，IGは，
\begin{equation}
 \mathrm{IG}_j(x;x')=(x_j-x'_j)\int_0^1
 \frac{\partial F(x'+\alpha(x-x'))}{\partial x_j}\,d\alpha
\end{equation}
で定義される\cite{sundararajan2017}．モデル出力が熱流束である場合，説明対象はその出力差である．分類の正解クラスの寄与と同じ意味へ置き換えない．出力をスケール変換する場合は，寄与値の単位も確認する．

基準入力と積分近似の選択を明記し，寄与和と$F(x)-F(x')$の差を確認する．異なるモデルで同じ色が付いても，基準出力やスケールが異なるため，色の濃さだけから重要度を直接比較しない．収束していても，物理的に実現しない積分経路や基準に依存した説明である可能性は残る．

\need{L03}{IGを主図へ追加する場合は，対象run・学習済みモデル・入力ID・基準・収束診断を揃え，入力対照との対応を記載する．現在の草稿は未照合の画像を主証拠にしないため，追加画像なしでも中心結論は記述できる．}

\section{実験装置の補足資料を追加する方針}
原実験資料の確認後，装置の模式図，センサ位置，電気計測からの熱流束算出，測定の不確かさを補足できる．引用または転載の場合は，原資料の著者，図番号，対象条件と必要な利用条件を明記する．現段階では，学部論文や先輩の装置図を対象日の原装置と推測して埋めない．本文の出典・ラベル・担当範囲は第2章に残し，付録への移動だけで出典確認を完了した扱いにはしない．
